% Part: second-order-logic
% Chapter: syntax-and-semantics
% Section: terms-formulas

\documentclass[../../../include/open-logic-section]{subfiles}

\begin{document}

\olfileid{sol}{syn}{frm}

\olsection{पदे आणि \printtoken{P}{formula}}

प्रथम-क्रम
तर्कशास्त्राप्रमाणेच
द्वितीय-क्रम
तर्कशास्त्रातील
अभिव्यक्तींची
रचना मूलभूत
शब्दसंग्रहापासून
होते. त्यात
\emph{!!{variable}s},
\emph{!!{constant}s},
\emph{!!{predicate}s}
आणि कधीकधी
\emph{!!{function}s}
असतात.
या चिन्हांपासून,
तार्किक संयोजके,
संख्यक आणि कंस
व स्वल्पविराम
यांसारखी
विरामचिन्हे
यांच्या साहाय्याने,
\emph{पदे} आणि
\emph{!!{formula}s}
तयार होतात.
फरक असा की
वस्तूंसाठीच्या
चलांव्यतिरिक्त
द्वितीय-क्रम
तर्कशास्त्रात
संबंधांसाठी
आणि फलनांसाठीही
चले असतात;
त्यांवर
संख्यकीकरण
करण्याचीही
मुभा असते.
म्हणून द्वितीय-क्रम
तर्कशास्त्राची
तार्किक चिन्हे
म्हणजे प्रथम-क्रम
तर्कशास्त्राची
चिन्हे आणि
यांखेरीज पुढील:

\begin{enumerate}
\item प्रत्येक
  स्थानसंख्या~$n$
  साठी द्वितीय-क्रम
  संबंध-!!{variable}s
  चा !!{denumerable}s
  संच:
  $\Obj V_0^n$, $\Obj V_1^n$, $\Obj V_2^n$, \dots
\item द्वितीय-क्रम
  फलन-!!{variable}s
  चा !!{denumerable}s
  संच:
  $\Obj u_0^n$, $\Obj u_1^n$, $\Obj u_2^n$, \dots
\end{enumerate}

प्रथम-क्रम
चल~$\Obj v_i$
यांसाठी आपण
$x$, $y$, $z$
ही अधिचले
वापरतो; त्याचप्रमाणे
$\Obj V_i^n$
यांसाठी
$X$, $Y$, $Z$
इत्यादी, आणि
$\Obj u_i^n$
यांसाठी
$u$, $v$
इत्यादी अधिचले
वापरू.

\begin{explain}
द्वितीय-क्रम
भाषेतील अतार्किक
चिन्हे प्रथम-क्रम
भाषेप्रमाणेच
निर्दिष्ट होतात:
तिची !!{constant}s,
!!{function}s
आणि !!{predicate}s
यांची यादी देऊन.

प्रथम-क्रम
तर्कशास्त्रात
!!{identity}~$\eq$
सहसा समाविष्ट
असते.
प्रथम-क्रम
तर्कशास्त्रात
भाषा~$\Lang{L}$
मधील अतार्किक
चिन्हे काही
लक्षणीय गोष्ट
व्यक्त करण्यासाठी
अत्यावश्यक असतात.
अतार्किक चिन्हे
न वापरणारी
!!{sentence}s
अर्थात असतात;
पण फक्त~$\eq$
असल्यास काही
लक्षणीय गोष्ट
सांगणे कठीण
असते.
द्वितीय-क्रम
तर्कशास्त्रात
आपल्याकडे संबंध-
आणि फलन-चलांचा
अमर्याद पुरवठा
असल्यामुळे,
अतार्किक चिन्हांचा
विशेष पुरवठा
नसतानाही
प्रथम-क्रम
भाषेत जे काही
सांगता येते
ते सांगता येते.
\end{explain}

\begin{defn}[द्वितीय-क्रम पदे]
भाषा~$\Lang L$
मधील \emph{द्वितीय-क्रम
पदांचा} संच
$\TrmSOL[L]$
हा \olref[fol][syn][frm]{defn:terms}
मध्ये पुढील
खंड वाढवून
परिभाषित होतो:
\begin{enumerate}
\item $u$ हे
  $n$-स्थानी
  फलन-चल
  आणि $t_1$,
  \dots, $t_n$
  ही पदे असतील,
  तर $\Atom{u}{t_1, \ldots, t_n}$
  हे पद असते.
\end{enumerate}
\end{defn}

\begin{explain}
म्हणून द्वितीय-क्रम
पद प्रथम-क्रम
पदासारखेच दिसते;
फक्त प्रथम-क्रम
पदात जिथे
!!a{function}~$\Obj{f^n_i}$
असते, तिथे
द्वितीय-क्रम
पदात त्याच्या
जागी फलन-चल~$\Obj{u^n_i}$
असू शकते.
\end{explain}

\begin{defn}[द्वितीय-क्रम \usetoken{s}{formula}]
भाषा~$\Lang L$
मधील \emph{द्वितीय-क्रम
!!{formula}s} चा
संच~$\FrmSOL[L]$
हा \olref[fol][syn][frm]{defn:formulas}
मध्ये पुढील
खंड वाढवून
परिभाषित होतो:
\begin{enumerate}
\item $X$ हे
  $n$-स्थानी
  विधेय-चल
  आणि $t_1$,
  \dots, $t_n$
  ही~$\Lang L$
मधील द्वितीय-क्रम
पदे असतील, तर
  $\Atom{X}{t_1,\ldots, t_n}$
  हे अणुरूप
  !!{formula}
  असते.

\tagitem{prvAll}{$!A$ हे !!a{formula} आणि $u$ हे फलन-चल असेल,
  तर $\lforall[u][!A]$ हे !!a{formula} असते.}{}

\tagitem{prvAll}{$!A$ हे !!a{formula} आणि $X$ हे विधेय-चल असेल,
  तर $\lforall[X][!A]$ हे !!a{formula} असते.}{}

\tagitem{prvEx}{$!A$ हे !!a{formula} आणि $u$ हे फलन-चल असेल,
  तर $\lexists[u][!A]$ हे !!a{formula} असते.}{}

\tagitem{prvEx}{$!A$ हे !!a{formula} आणि $X$ हे विधेय-चल असेल,
  तर $\lexists[X][!A]$ हे !!a{formula} असते.}{}
\end{enumerate}
\end{defn}

\end{document}
